\chapter{基于全双工建模的自然语音交互大模型}\label{chap:full-duplex}

\section{引言}

前两章分别从响应延迟和语音生成质量两个维度推进了语音交互模型的发展。第\ref{chap:llama-omni}章提出的LLaMA-Omni通过非自回归语音解码器实现了低至236毫秒的响应延迟，第\ref{chap:llama-omni2}章提出的LLaMA-Omni 2进一步引入自回归流式语音解码器，在保持实时交互的同时大幅提升了语音生成的准确性和自然度。然而，上述两个模型均面向基于轮次的对话场景，即模型接收一段完整的语音输入，生成对应的文本和语音响应。在实际部署中，这类模型通常需要配合前置的语音活动检测（Voice Activity Detection, VAD）模块来判断用户何时结束发言，再将检测到的语音片段输入模型以获取响应。尽管这种级联方案在大多数场景下能够正常工作，但存在两方面的固有局限。

一方面，系统的整体表现受限于VAD模块的检测精度。VAD的误检会导致语音被过早切分，使模型基于不完整的问题开始回答；而VAD的漏检则会延迟语音切分的时机，增加用户感知到的响应延迟。由于VAD通常仅依赖声学特征的变化来判断语音活动的起止，缺乏对语义信息的理解，这些错误在实际应用中难以完全避免。另一方面，基于轮次的交互模式限制了更灵活的对话能力。在真实的双人对话中，除了一问一答的基本模式外，还存在多种复杂的交互场景：用户在发言过程中可能出现思考停顿，此时模型不应将其误判为发言结束而仓促回答，而应等待用户继续；用户也可能在模型回复的过程中主动打断，模型需要及时停止当前回复并处理新的输入。基于VAD的级联方案将回复时机的决策完全交给一个小型前端模块，无法根据对话的语义上下文做出合理判断，这从根本上限制了系统的交互能力上限。

针对上述问题，全双工语音交互模型开始受到越来越多的关注。与基于轮次的模型不同，全双工模型在运行时持续接收环境中的音频输入，由模型自主决定何时开始说话、何时停止说话。这意味着各种对话时机的决策从外部的VAD模块转移到大语言模型内部，与对话内容的理解和生成进行端到端的联合学习，从而具备更高的交互能力上限。

然而，全双工模型的构建面临多方面的技术挑战。其一，与基于轮次的模型在每个时刻只需处理用户语音或模型语音不同，全双工模型要求在每个时刻同时处理输入和输出两个语音信道，这与语言模型单向顺序生成的基本建模方式相悖，如何使模型具备同时处理多个信道序列的能力是首要挑战。其二，为保证生成回复的质量，通常需要像基于轮次的模型一样引入额外的文本信道来引导语音生成，但文本以子词为基本粒度，而语音以固定时间步为基本粒度，两者之间存在粒度不匹配的问题，如何在全双工框架下同时协调文本信道与两个语音信道需要专门设计。其三，真实的全双工对话数据极为稀缺，现有的大规模语音数据多为基于轮次的对话格式，如何充分利用这些数据辅助全双工模型的训练，或如何基于现有的基于轮次的语音-语言模型扩展到全双工场景，是缓解数据瓶颈的关键。其四，部分现有工作通过添加额外的分类头等模块来决定模型的输出时机，这种方式在复杂的全双工场景下泛化能力有限，如何设计更具通用性的建模方式也是一个重要问题。

为解决上述挑战，本章提出了一种基于多信道交织的全双工语音交互模型。该模型将用户语音、模型回复文本和模型回复语音三个信道按照固定比例组织成交织序列，统一由大语言模型进行建模。用户语音和模型语音均使用GLM-4-Voice\citep{glm4voice}的语音分词器进行离散化，帧率为每秒12.5个词元。文本信道包含模型回复的文本内容以及若干特殊标记：\texttt{[SILENCE]}表示模型当前应保持静默，\texttt{[START]}表示模型开始回复，\texttt{[PAD]}表示文本生成已结束但对应语音仍在生成，\texttt{[EPAD]}表示语音生成结束。由于三个信道按固定比例交织，文本信道在时间轴上与语音信道完全对齐，每个时刻的语音词元在文本信道均有对应的标记来指示当前的对话状态，从而解决了文本与语音建模粒度不匹配的问题。在训练策略上，本章基于GLM-4-Voice进行轻量的全双工微调。GLM-4-Voice已经在百万小时级语音数据上完成预训练，并经过基于轮次的语音对话微调，具备了较强的语音理解和生成能力。通过少量全双工数据的微调，模型即可适应全双工的建模范式，同时充分利用预训练和对话微调阶段已习得的知识。此外，这种多信道交织的建模方式无需引入额外的分类头等决策模块，面对不同的全双工场景均可通过统一的文本信道标记规则进行处理，具备更好的泛化能力。为进一步提升模型在回复时机上的决策能力，本章还引入了直接偏好优化（Direct Preference Optimization, DPO）\citep{rafailov2023dpo}，通过降低不期望行为的生成概率来强化模型的交互表现。

在轮替和打断两类全双工场景上的实验结果表明，仅使用200000条数据进行全双工微调，模型即可具备自主判断回复时机的能力，轮替成功率达到93\%，打断成功率达到XX\%，同时对话质量相比Moshi\citep{kyutai2024moshi}等基线模型取得了显著提升。
